
\section{The support law}\label{sec:supportlaw}

This section proves Theorem~\ref{main:supportlaw}. Throughout, a primitive
positive Parker configuration is written with the three square progressions
$(r_i^2,s_i^2,t_i^2)$ of common difference $d$, and
\[
u_i=t_i-r_i,\qquad v_i=t_i+r_i,\qquad z_i=u_i/v_i\in(0,1),
\qquad u_iv_i=2d .
\]

\subsection{The torus identity and nondegeneracy}

Substituting $u_i,v_i$ into the interaction $s_1^2+s_2^2=2s_0^2$ and using
$4s_i^2=u_i^2+v_i^2$ gives, after division by $2d=u_iv_i$, exactly
\[
z_1+z_1^{-1}+z_2+z_2^{-1}=2\,(z_0+z_0^{-1}) :
\tag{L.1}
\]
multiplying by $z_0/2$, the five-term unit equation
\[
\tfrac{z_0z_1}{2}+\tfrac{z_0}{2z_1}+\tfrac{z_0z_2}{2}+\tfrac{z_0}{2z_2}
-z_0^{2}=1
\tag{L.1$'$}
\]
in the multiplicative group generated by the $z_i$.

\emph{Nondegeneracy.} On the Parker open --- all $z_i\in(0,1)$,
pairwise-distinct progressions --- no nonempty proper subsum of the left
side of (L.1$'$) vanishes. Two exact facts drive every case. First, all
$z_i$ and $z_i^{-1}$ lie in the squareclass of $2d$: indeed
$z_i\cdot 2d=u_i^{2}$ and $2d/z_i=v_i^{2}$, so any relation of the shape
$z_i=2z_j$, $z_iz_j=2$, $z_i=2z_j^{-1}$ or $2s_i^2=u_0^2$ would make $2$ a
rational square. Second, the interaction itself. A vanishing subsum
containing none of the four positive terms is empty; one containing only
positive terms is positive. The remaining cases pair $-z_0^{2}$ with one,
two, or three positive terms. With \emph{one}: $z_0z_1/2=z_0^{2}$ gives
$z_1=2z_0$, hence $u_1^{2}=2u_0^{2}$ --- excluded, and likewise for the
other three singletons. With \emph{three}: the subsum equals
$1-(\text{omitted term})$ by (L.1$'$) itself, so the omitted term equals
$1$, giving $z_0z_1=2$ or one of its three analogues --- excluded the same
way. With \emph{two} same-index terms:
$z_1+z_1^{-1}=2z_0$ gives $4s_1^{2}/(2d)=u_0^{2}/d$, i.e.\
$2s_1^{2}=u_0^{2}$ --- excluded. With \emph{two} pure cross terms,
$z_1+z_2=2z_0$: combined with the interaction
$(u_1^2+v_1^2)+(u_2^2+v_2^2)=2(u_0^2+v_0^2)$ and $v_i=2d/u_i$, the two
relations force, for $x=(u_0/u_1)^{2}$ and $y=(u_0/u_2)^{2}$, both
$x+y=2$ and $xy=1$, hence $x=y=1$ and $z_1=z_2=z_0$ --- against pairwise
distinctness (and symmetrically for $z_1^{-1}+z_2^{-1}=2z_0$). With
\emph{two} mixed cross terms, $z_1+z_2^{-1}=2z_0$: the same elimination
collapses to $(u_0^{2}u_2^{2}-4d^{2})^{2}=0$, i.e.\ $u_0=v_2$ and
$v_0=u_2$, so $z_0z_2=1$ --- impossible with both in $(0,1)$. Hence
(L.1$'$) is nondegenerate in the sense required by the subspace theorem.

\subsection{The fixed-support count}

Fix a finite prime set $\mathcal S\ni2$ and restrict to configurations with
$\supp(2d)\subseteq\mathcal S$. Modulo common rational-square scaling of the
entries, the triple $(z_0,z_1,z_2)$ determines the configuration, and each
$z_i$ lies in the group of $\mathcal S$-units of rank at most
$3|\mathcal S|$ relevant to (L.1). Theorem~1.1 of Evertse, Schlickewei and
Schmidt~\cite{ESS02} bounds the nondegenerate solutions of a five-term unit
equation over a group of rank $r$ by $\exp\bigl(30^{15}(r+1)\bigr)$;
applying it with $r\le3|\mathcal S|$ gives at most
\[
\exp\bigl(30^{15}(3|\mathcal S|+1)\bigr)
\]
root-primitive configurations. The bound counts; it does \textbf{not}
enumerate, and no effective procedure follows from it.

\subsection{Tropical and Legendre laws}

For odd $p$ the six valuations occurring in (L.1) are
$\pm\nu_i$ with $\nu_i=v_p(z_i)$; in a vanishing non-Archimedean sum the
least valuation is attained at least twice, which proves that the largest
$|\nu_i|$ occurs at least twice. At $p=2$ the same argument applies to
$\min\{\nu_1,-\nu_1,\nu_2,-\nu_2,1+\nu_0,1-\nu_0\}$. Writing $e=v_p(2d)$
and $a_i=v_p(u_i)$, the relation $u_iv_i=2d$ gives $\nu_i=2a_i-e$ with
$0\le a_i\le e$; call an index an \emph{endpoint} if $a_i\in\{0,e\}$ and
\emph{interior} if $0<a_i<e$. For odd $p$, an interior index puts $p$ in
every root of its progression: $p$ divides both $u_i=t_i-r_i$ and
$v_i=t_i+r_i$, hence $t_i$ and $r_i$, and then $4s_i^2=u_i^2+v_i^2$ forces
$p\mid s_i$. Root primitivity therefore forbids three interior indices, so
at least one index is an endpoint; and an endpoint attains the maximal
value $|\nu_i|=e$, so the tropical tie forces a second endpoint:
\textbf{at least two indices are endpoints at every odd support prime}.
If exactly two indices are endpoints, comparing
the unit squareclasses of the two least-valuation terms forces $-1$ to be a
quadratic residue for the pair $\{1,2\}$ and $2$ for a pair containing
$0$; hence the exact endpoint pairs obey
\[
\begin{array}{r|l}
p\bmod 8 & \text{allowed exact endpoint pairs}\\\hline
1 & \{0,1\},\ \{0,2\},\ \{1,2\}\\
3 & \text{none; all three indices are endpoints}\\
5 & \{1,2\}\ \text{only}\\
7 & \{0,1\},\ \{0,2\}\ \text{only}
\end{array}
\]
These are necessary local laws, not sufficient conditions.

\subsection{Exclusion of support $\{2,3\}$}

By Lemma~1 of~\cite{PTZ15} every entry of a primitive configuration is
$1\bmod24$, so all roots are coprime to $6$ and $24\mid d$. Assume
$\supp(2d)\subseteq\{2,3\}$ and put $e_p=v_p(2d)$. At $2$, both $u_i,v_i$
are even and exactly one has valuation $1$, so $v_2(z_i)=\pm(e_2-2)$; at
$3$, exactly one of them is divisible by $3$, so $v_3(z_i)=\pm e_3$; there
are no other primes. Hence $z_i\in\{2^{\pm(e_2-2)}3^{\pm e_3}\}$, and since
$z\mapsto z^{-1}$ pairs the four sign choices, $f(z)=z+z^{-1}$ takes at
most two values on the $z_i$. Three elements of a two-element set can
satisfy $f(z_1)+f(z_2)=2f(z_0)$ only if all three values agree, and $f$ is
strictly decreasing on $(0,1)$, so $z_0=z_1=z_2$; with the same $K=2d$ the
formulas $u_i=\sqrt{Kz_i}$, $v_i=\sqrt{K/z_i}$ then force the three
progressions to coincide, contradicting pairwise distinctness. This
completes the proof of Theorem~\ref{main:supportlaw}.

\begin{remark}[Scope]
Everything here is a necessary-condition layer on the common difference.
The subspace-theorem count is imported~\cite{ESS02}; the congruence input
is~\cite{PTZ15}; the tropical mechanism is elementary. No effectivity, no
global support-size bound and no all-support emptiness follows, and none is
claimed.
\end{remark}
